\begin{table}[!tbp]
\caption{Does ordering by slowness \emph{find} the persistent factor? All four score
columns are the persistent factor read by a linear probe from a 16-dimensional subspace
of the \emph{same mechanism-free} embedding (macro-F1 on real data, accuracy on
synthetic); only the choice of subspace differs. $C$ is the fraction of the reachable
range $(\text{full}-\text{random})$ that the slowness ordering recovers; $n=5$ seeds,
$\pm$ is the seed standard deviation. \textbf{Read the sign, not the size}: the
denominator is small, so the ratio is noisy, but real data and synthetic data do not
overlap.}
\label{tab:metricC}
\centering
\small
\begin{tabular}{llcccc}
\hline
Dataset & Stem & SFA slow half & Random 16 & Full 64 & $C$ \\
\hline
Synthetic & HGLP-Reg & 0.862 & 0.856 & 0.886 & $\mathbf{-0.22} \pm 0.72$ \\
Synthetic & HGLP-NCE & 0.986 & 0.986 & 0.985 & $\mathbf{-1.49} \pm 1.24$ \\
PTB-XL & HGLP-Reg & 0.770 & 0.750 & 0.778 & $\mathbf{+0.42} \pm 0.73$ \\
PTB-XL & HGLP-NCE & 0.791 & 0.776 & 0.793 & $\mathbf{+0.78} \pm 0.61$ \\
HAPT & HGLP-Reg & 0.898 & 0.881 & 0.920 & $\mathbf{+0.35} \pm 0.59$ \\
HAPT & HGLP-NCE & 0.909 & 0.876 & 0.929 & $\mathbf{+0.56} \pm 0.31$ \\
\hline
\end{tabular}
\end{table}
